\section{Language Electives}\label{language-electives}

\stanceinfo{\textbf{CFES Congress 2018} [2018-01-07]}{2024-01-02}

\officialstance{The Canadian Federation of Engineering Students believes that Canadian engineering students should have the opportunity to take elective language skills as part of their undergraduate program to enrich their engineering degree.}

\stanceheading{The Students' Position}
\begin{itemize}
 \item In 2018, the Canadian Engineering Accreditation Board (CEAB) updated their \href{https://engineerscanada.ca/sites/default/files/accreditation/2021-2022-cycle/accreditation-criteria-procedures-2020.pdf}{Accreditation Criteria \& Procedures} [2] to explicitly accept language electives as part of their ``Complementary Studies'' requirement for engineering programs across Canada; however, around 23\% of accredited engineering schools in Canada have yet to reflect this development within their faculties.
 \item In 2022, the criterion was removed 3.4.5.2 as listed below, where the CFES has found it remains listed on certain complementary studies university websites:
 \begin{itemize}
  \item [REMOVED] Criterion 3.4.5.2 Language instruction may be included within complementary studies provided it is not taken to fulfill an admission requirement. Furthermore, curriculum content that principally imparts language skills can be counted toward the required Accreditation Unit (AU) of complementary studies but cannot be used to satisfy the requirements for subject matter that deals with central issues, methodologies, and thought processes of the humanities and social sciences.
  \item CFES agrees with the removed statement, as even subject matter that deals with other languages is inclusive to the enrichment of engineering students apart from their technical courses.
 \end{itemize}
 \item Knowledge of an additional language and surrounding subject matter is a valuable asset for the future of an engineering student and students must be given a full and appropriate range of language elective options.
\end{itemize}

\stanceheading{Recommendations to Partners, Stakeholders, and Other Entities}
\begin{itemize}
 \item The CFES calls on the CEAB to assist the final 23\% of Canadian institutions with the interpretation of their existing policy on complementary electives by publishing an official interpretive statement on the subject outlining additional criteria.
 \item The CFES calls on Engineering Deans Canada (EDC) to advocate for its member faculties to review and reconsider their elective offerings, and to allow engineering students to enroll in language courses that count towards their degree requirements, outside and including both Canadian official languages.
 \item The CFES calls on EDC to update their policies and information available on their academic calendars to accurately describe and follow the CEAB's Policies and Procedures as to what counts as a complementary course.
 \item The CFES calls on engineering student societies to support national advocacy efforts and leverage the CEAB policy update to advocate for the implementation of language electives in their program to reap the benefits of multilingualism in a diverse country and an increasingly globalized world.
 \item The CFES calls on their members to question the curriculum currently in place, and pursue studies that inspire them in language studies that would enrich their engineering education even if additional permissions are necessary to pursue.
\end{itemize}

\stanceheading{What the CFES is doing}
\begin{itemize}
 \item The CFES' discussion on the interpretation of language course policy over two mandate periods was accepted by the CEAB's Policies and Procedures Committee who approved language sections as open electives before President's Meeting 2019.
 \item The CFES has created and continues to update an advocacy package to guide students in advocating for language electives at an institutional level.
 \item The CFES has studied the language elective options within accredited engineering universities in Canada, compiling a list of universities currently not following policy.
 \item The CFES has presented data on CEAB universities that do not currently follow policy, or have outdated information to begin that conversation between EDC and CEAB.
\end{itemize}

\stanceheading{What the CFES plans to do}
\begin{itemize}
 \item The CFES will provide resources on language electives to equip its member schools with the information they require to advocate for elective changes at individual institutions.
 \item The CFES will provide an analysis on the benefits of having language electives available as complementary studies and how it enriches the education of an engineer.
 \item The CFES will work with EDC to make them aware of policies and documentation that go against or are not completely accurate based on current CEAB policy and procedure.
 \item The CFES will build on past advocacy efforts and continue the discussion on the importance of bilingualism/multilingualism during external partner presentations on the future of engineering education.
\end{itemize}

\newpage
